\section{Directional and curved finite parts of double Hurwitz zeta}
\label{sec:direction}

The incoming report asks what changes when the exponent perturbations
are independent, and whether its diagonal finite parts depend on the
approach path. At depth two these questions have a particularly explicit
answer. The meromorphic continuation itself belongs to the established
multiple-zeta theory \cite{mow}. The result below identifies its full
polar germ in the present normalization, then computes the resulting
path dependence and all regular Taylor coefficients.

For $a>0$ define
\begin{equation}\label{eq:doublehurwitz}
 T(s,t;a)=\sum_{n>m\ge0}(n+a)^{-s}(m+a)^{-t}.
\end{equation}
The defining series is absolutely convergent when
$\Rea s>1$ and $\Rea(s+t)>2$. Its order convention is essential: the
outer, larger index carries $s$. Put
\[
 g_a(u)=\zeta(1+u,a)-\frac1u
       =\sum_{j\ge0}g_j(a)u^j,
 \qquad g_j(a)=\frac{(-1)^j\gamma_j(a)}{j!}.
\]

\begin{theorem}[Exact polar germ at $(1,1)$]\label{thm:doublegerm}
There is a function $H_a(u,v)$ holomorphic near $(0,0)$ such that
\begin{equation}\label{eq:doublegerm}
 \boxed{T(1+u,1+v;a)
   =\frac1{u(u+v)}+\frac{g_a(v)}u+H_a(u,v).}
\end{equation}
It satisfies
\begin{align}
 H_a(u,v)+H_a(v,u)
   &=g_a(u)g_a(v)-\zeta(2+u+v,a),\label{eq:Hsym}\\
 H_a(0,0)&=C_a:=\frac{\gamma_0(a)^2-\zeta(2,a)}2.
 \label{eq:H00}
\end{align}
Thus the two local polar divisors are $u=0$ and $u+v=0$; in particular
there is no separate polar divisor $v=0$.
\end{theorem}

\begin{proof}
Write $x=n+a$, and use the inner finite sum to obtain, first in the
absolute-convergence domain,
\[
 T(1+u,1+v;a)=\zeta(1+u,a)\zeta(1+v,a)
       -\sum_{n\ge0}x^{-1-u}\zeta(1+v,x).
\]
The pole-subtracted Euler--Maclaurin remainder
\[
 E(v,x)=\zeta(1+v,x)-\frac{x^{-v}}v-\frac12x^{-1-v}
\]
is holomorphic in $v$ near zero and is
$O(x^{-2-\Rea v})$ uniformly on compact $v$-sets there. Hence
\[
 R_a(u,v)=\sum_{n\ge0}(n+a)^{-1-u}E(v,n+a)
\]
converges normally in a neighborhood of $(0,0)$, as do its fixed
derivatives. We obtain
\begin{equation}\label{eq:doubleEM}
 T=\zeta(1+u,a)\zeta(1+v,a)
   -\frac{\zeta(1+u+v,a)}v
   -\frac12\zeta(2+u+v,a)-R_a(u,v).
\end{equation}
The apparent singularity at $v=0$ cancels. Substitution of
$\zeta(1+u,a)=u^{-1}+g_a(u)$ gives \eqref{eq:doublegerm}, with
\begin{equation}\label{eq:Hrepresentation}
 H_a(u,v)=g_a(u)g_a(v)
  +\frac{g_a(u)-g_a(u+v)}v
  -\frac12\zeta(2+u+v,a)-R_a(u,v).
\end{equation}
The divided difference is holomorphic, since its numerator vanishes at
$v=0$. Finally, the finite-sum stuffle identity continues to
\[
 T(s,t;a)+T(t,s;a)=\zeta(s,a)\zeta(t,a)-\zeta(s+t,a).
\]
Insert \eqref{eq:doublegerm} to get \eqref{eq:Hsym} and then
\eqref{eq:H00}. All steps are identities of meromorphic germs, so no
choice of a one-dimensional limiting path is hidden in the proof.
\end{proof}

\subsection{Straight directions and their exact discrepancy}

Here $\FP_{\eps=0}$ means the coefficient of $\eps^0$ in the Laurent
expansion after the specified substitution. It is not a joint value at
$(u,v)=(0,0)$.

\begin{corollary}[Every transverse straight direction]\label{cor:direction}
Let $c,d\in\C$, with $c\ne0$ and $c+d\ne0$. Then
\begin{align}\label{eq:direction-expansion}
 T(1+c\eps,1+d\eps;a)
  =\frac1{c(c+d)\eps^2}
     +\frac{\gamma_0(a)}{c\eps}
     +C_a-\frac dc\gamma_1(a)+O(\eps),
\end{align}
and in particular
\begin{equation}\label{eq:direction-FP}
 \boxed{\FP_{\eps=0}T(1+c\eps,1+d\eps;a)
   =\frac{\gamma_0(a)^2-\zeta(2,a)}2-\frac dc\gamma_1(a).}
\end{equation}
The case $d=0$ is included. The difference between two admissible
directions is exactly
$-\gamma_1(a)(d/c-d'/c')$.
\end{corollary}
\begin{proof}
Substitute in \eqref{eq:doublegerm}, expand $g_a(d\eps)$ through its
linear term, and use \eqref{eq:H00}.
\end{proof}

For $c=d=1$, this agrees with the diagonal identity
$T(s,s;a)=(\zeta(s,a)^2-\zeta(2s,a))/2$.
For $d=0$, the finite part is $C_a$ instead. A curve contained in
$u=0$ or $u+v=0$ cannot be evaluated by restriction of the displayed
meromorphic function without an additional prescription. Curves tangent
to those divisors can have higher pole orders and are outside the
transversality statement.

\subsection{Curvature and reparametrization also matter}

The tangent direction alone does not determine the constant Laurent
coefficient. The double pole reads the cubic jet of a curve.

\begin{theorem}[Finite part along an analytic transverse curve]
\label{thm:curved}
Suppose
\[
 u=c_1\eps+c_2\eps^2+c_3\eps^3+O(\eps^4),\qquad
 v=d_1\eps+d_2\eps^2+d_3\eps^3+O(\eps^4),
\]
where $c_1\ne0$ and $A:=c_1+d_1\ne0$. Set
$B=c_2+d_2$ and $C=c_3+d_3$. Then
\begin{align}\label{eq:curved-FP}
 \FP_{\eps=0}T(1+u,1+v;a)
 ={}&C_a-\gamma_1(a)\frac{d_1}{c_1}
             -\gamma_0(a)\frac{c_2}{c_1^2}\notag\\
 &+\frac1{c_1A}\left\{
       \frac{c_2^2}{c_1^2}+\frac{B^2}{A^2}
       +\frac{c_2B}{c_1A}-\frac{c_3}{c_1}-\frac CA
                         \right\}.
\end{align}
\end{theorem}
\begin{proof}
The constant term of the holomorphic part is $C_a$. In $g_a(v)/u$,
the constant terms are $-\gamma_0(a)c_2/c_1^2$ and
$-\gamma_1(a)d_1/c_1$. Expanding both reciprocal factors in
$1/[u(u+v)]$ to second order after factoring out $c_1A\eps^2$
gives the braces. Terms of order four in the curve cannot contribute.
\end{proof}

For example, $u=\eps$, $v=\eps+\lambda\eps^2$ gives
\begin{equation}\label{eq:curve-example}
 \FP T(1+\eps,1+\eps+\lambda\eps^2;a)
       =C_a-\gamma_1(a)+\frac{\lambda^2}{8}.
\end{equation}
Thus even curves with the same tangent as the diagonal have different
finite parts. This is an explicit answer to the path-dependence question,
not merely an assertion that multivariate regularization can be ambiguous.
Even the same geometric diagonal, reparametrized by
$u=v=\eps+\lambda\eps^2+\mu\eps^3$, has finite part
\[
 C_a-\gamma_1(a)-\lambda\gamma_0(a)+\frac{3\lambda^2-2\mu}{2}.
\]
This also follows directly from the exact diagonal stuffle identity,
and independently checks the cubic-jet terms in \eqref{eq:curved-FP}.

\subsection{All regular Taylor coefficients in Stieltjes coordinates}

The convergent sums $S_{m,n}(a)$ from \eqref{eq:Sdef} give a direct
coordinate system for the complete holomorphic remainder.

\begin{theorem}[The complete two-variable regular jet]\label{thm:Hcoeff}
Write $H_a(u,v)=\sum_{m,n\ge0}h_{m,n}(a)u^mv^n$. For all $m,n\ge0$,
\begin{align}\label{eq:Hcoeff}
 h_{m,n}(a)
 ={}&\frac{(-1)^{m+n}}{m!n!}
 \left\{\gamma_m(a)\gamma_n(a)
          +\frac{\gamma_{m+n+1}(a)}{n+1}-S_{m,n}(a)\right\}\notag\\
 &-\frac{\zeta^{(m+n)}(2,a)}{2m!n!}.
\end{align}
For a transverse straight direction, every regular coefficient is
therefore explicit:
\begin{equation}\label{eq:direction-all}
 [\eps^q]T(1+c\eps,1+d\eps;a)
 =\frac{(-1)^{q+1}\gamma_{q+1}(a)d^{q+1}}{c(q+1)!}
       +\sum_{m+n=q}h_{m,n}(a)c^md^n,
 \qquad q\ge0.
\end{equation}
\end{theorem}
\begin{proof}
Normal convergence permits coefficient extraction in $R_a$. The
coefficient of $v^n$ in $E(v,x)$ is
\[
 \frac{(-1)^n}{n!}\left\{
 \gamma_n(x)+\frac{\log^{n+1}x}{n+1}-\frac{\log^n x}{2x}\right\}.
\]
Consequently $[u^mv^n]R_a=(-1)^{m+n}S_{m,n}/(m!n!)$.
The coefficient of the divided difference in
\eqref{eq:Hrepresentation} is
$-\binom{m+n+1}{n+1}g_{m+n+1}(a)$. Combining these two observations
with the coefficients of $g_a(u)g_a(v)$ and
$\zeta(2+u+v,a)/2$ proves \eqref{eq:Hcoeff}. Substitution in
\eqref{eq:doublegerm} proves \eqref{eq:direction-all}.
\end{proof}

Taking symmetric coefficients in \eqref{eq:Hsym} and using
\eqref{eq:Hcoeff} recovers \eqref{eq:Stwo}. Thus the nonlinear
Stieltjes summation law has two independent proofs: discrete telescoping
and the meromorphic double-zeta stuffle identity. This compatibility is
also an effective sign and normalization check.

